\chapter{PTEN Genetic Testing}

\section*{What Is This Test?}
PTEN (phosphatase and tensin homolog) genetic testing analyzes the PTEN tumor suppressor gene located on chromosome 10q23.3 for germline pathogenic variants. The PTEN protein is a dual-specificity phosphatase that dephosphorylates both protein substrates and lipid second messengers, specifically phosphatidylinositol-3,4,5-trisphosphate (PIP3), the product of PI3 kinase. By dephosphorylating PIP3 to PIP2, PTEN directly antagonizes the PI3K-AKT-mTOR signaling pathway, which promotes cell growth, proliferation, and survival. Loss-of-function germline mutations in PTEN cause PTEN hamartoma tumor syndrome (PHTS), which includes Cowden syndrome (multiple hamartomas, macrocephaly, and markedly increased risks of breast cancer [25--50\% lifetime], thyroid cancer [10--35\% non-medullary], endometrial cancer [25--30\%], renal cell carcinoma, and colorectal cancer), Bannayan-Riley-Ruvalcaba syndrome (macrocephaly, intestinal hamartomatous polyposis, lipomas, and pigmented macules of the glans penis), and Lhermitte-Duclos disease (dysplastic cerebellar gangliocytoma). PTEN testing is also performed on tumor tissue in oncology because somatic PTEN loss occurs in numerous sporadic cancers (prostate, endometrial, glioblastoma, breast) and may predict responsiveness or resistance to specific therapies, including mTOR and PI3K inhibitors.

\section*{Alternative Names}
PTEN Mutation Analysis; PTEN Sequencing; Cowden Syndrome Testing; PTEN Hamartoma Tumor Syndrome Testing; PHTS Genetic Testing; MMAC1 Testing

\section*{Principle}
Germline PTEN testing is performed by next-generation sequencing (NGS) of the coding exons and flanking splice sites of PTEN, with deletion/duplication analysis by MLPA or chromosomal microarray to detect large rearrangements. Somatic PTEN testing is performed on formalin-fixed paraffin-embedded tumor tissue by NGS (comprehensive genomic profiling panel). PTEN testing is often included in multigene hereditary cancer panels. Immunohistochemistry (IHC) for PTEN protein expression in tumor tissue can serve as a screen for PTEN loss, but PTEN IHC is not standardized for clinical use and germline/Somatic sequencing is required for definitive diagnosis.

\section*{History}
PTEN was discovered simultaneously in 1997 by two independent groups: Parsons et al. at Columbia University identified PTEN on 10q23 as the tumor suppressor gene deleted in glioblastoma, prostate cancer, and breast cancer, while Li et al. identified the same gene, MMAC1 (mutated in multiple advanced cancers 1), on 10q23.3 as a tumor suppressor frequently mutated in advanced cancers. Germline PTEN mutations were linked to Cowden syndrome by Liaw et al. in 1997. The discovery of the PTEN-PI3K-AKT-mTOR signaling pathway revolutionized the understanding of cancer biology and led to the development of therapeutic mTOR inhibitors (sirolimus, everolimus, temsirolimus) and PI3K/AKT inhibitors for PTEN-null cancers.

\section*{Why This Test Is Ordered}
Diagnosis of PTEN hamartoma tumor syndrome (Cowden syndrome): patients with macrocephaly, multiple hamartomas, breast cancer, thyroid cancer, endometrial cancer, and/or pathognomonic mucocutaneous lesions (trichilemmomas, acral keratoses, oral papillomatous papules). NCCN guidelines for hereditary breast and ovarian cancer recommend PTEN testing in patients with personal history of breast cancer and macrocephaly. Somatic PTEN testing in tumor tissue for identification of actionable PI3K-AKT-mTOR pathway activation in the context of clinical trials or targeted therapy. PTEN loss by IHC combined with ERG rearrangements or PTEN deletion by FISH in prostate cancer for prognostication and clinical trial eligibility.

\section*{Is This a Primary or Secondary Test}
This is a primary diagnostic test for PTEN hamartoma tumor syndrome when clinical criteria are met.

\section*{How This Test Is Performed}
Blood sample in EDTA (lavender-top) tube for germline testing. Tumor tissue (FFPE block) for somatic testing.

\section*{What Preparation Is Needed}
Pre-test genetic counseling is recommended. No fasting or medication restrictions.

\section*{Contraindications and Precautions}
The primary contraindication is the absence of informed consent and pre-test genetic counseling, given the profound implications of a positive result for the patient and family. Precautions: PTEN testing is informative mainly in individuals who meet clinical criteria for PTEN hamartoma tumor syndrome; testing unaffected individuals without a suggestive personal or family history has a low yield. A variant of uncertain significance (VUS) is not actionable and should not guide management; VUS classification may change over time as data accumulate. Because PTEN-associated cancers (breast, thyroid, endometrial) overlap with other hereditary syndromes, a negative PTEN result does not exclude an inherited predisposition, and multigene panel testing should be considered when clinical features are suggestive. Somatic PTEN findings in tumor tissue reflect the tumor only and do not establish a germline syndrome.

\section*{Risks}
Minimal physical risk (bruising, pain at the venipuncture site). The greater risks are psychological (anxiety, depression) related to a positive result and its implications for family members. The Genetic Information Nondiscrimination Act (GINA) of 2008 protects against genetic discrimination in health insurance and employment but does not apply to life, disability, or long-term care insurance.

\section*{What Happens After the Test}
Results typically return in 2--4 weeks and are reported as positive (pathogenic or likely pathogenic variant), negative, or VUS. A positive result confirms PTEN hamartoma tumor syndrome and triggers NCCN surveillance, including annual breast MRI and mammography starting at age 30--35, thyroid ultrasound, renal imaging, colonoscopy every 5 years starting at age 35, annual endometrial biopsy by age 30--35 with discussion of risk-reducing hysterectomy, and annual dermatologic examination. Post-test genetic counseling is provided, and cascade testing of at-risk relatives (each with a 50\% probability of carrying the variant) is recommended.

\section*{Understanding Results}

\subsection*{Pathogenic/Likely Pathogenic Variant Detected (Positive)}
Confirms PTEN hamartoma tumor syndrome. NCCN screening guidelines: (1) Breast cancer: annual mammogram and breast MRI with contrast starting at age 30--35 (or 10 years before the earliest breast cancer in the family). Discussion of risk-reducing mastectomy. (2) Endometrial cancer: discussion of risk-reducing hysterectomy after completion of childbearing. Annual endometrial biopsy by age 30--35. (3) Thyroid cancer: annual thyroid ultrasound starting at diagnosis (pediatric) or at age 18. (4) Renal cell carcinoma: renal ultrasound or CT/MRI every 1--2 years starting at age 40. (5) Colorectal cancer: colonoscopy every 5 years starting at age 35. (6) Dermatologic examination annually. (7) Neuropsychological evaluation for autism spectrum disorder in children.

\subsection*{No Pathogenic Variant Detected (Negative)}
No evidence of PTEN hamartoma tumor syndrome. Management based on personal and family history. Other hereditary cancer syndromes should be considered if clinical features are suggestive.

\subsection*{Variant of Uncertain Significance (VUS)}
Not actionable clinically. Periodic re-evaluation of variant classification as new data emerge.

\section*{Normal Results}
No pathogenic variant detected. Constitutional PTEN pathogenic variants in the general population are extremely rare.

\section*{Follow-up Tests}
Cascade testing of at-risk family members. Cancer surveillance as per NCCN guidelines (annual breast MRI/mammogram, thyroid ultrasound, colonoscopy, renal imaging, endometrial assessment). Somatic PTEN testing in tumor tissue (comprehensive genomic profiling). PTEN IHC in tumor tissue as a screen for PTEN loss.

\section*{References}
\begin{enumerate}
  \item Pilarski R, Burt R, Kohlman W, et al. Cowden syndrome and the PTEN hamartoma tumor syndrome. J Natl Compr Canc Netw. 2013;11(8):1000--1010.
  \item Li J, Yen C, Liaw D, et al. PTEN, a putative protein tyrosine phosphatase gene mutated in human brain, breast, and prostate cancer. Science. 1997;275(5308):1943--1947.
  \item NCCN Genetic/Familial High-Risk Assessment: Breast, Ovarian, and Pancreatic, Version 3.2024.
  \item Tan MH, Mester JL, Ngeow J, et al. Lifetime cancer risks in individuals with germline PTEN mutations. Clin Cancer Res. 2012;18(2):400--407.
  \item Marsh DJ, Dahia PL, Zheng Z, et al. Germline mutations in PTEN are present in Bannayan-Zonana syndrome. Nat Genet. 1997;16(4):333--334.
\end{enumerate}

\clearpage
